\section{Systemic Interaction Benchmarks}

\subsection{The GST Relation: Geometric Basis}
The Gatto-Sartori-Tonin (GST) relation conceptually links the gauge sector to the flavor sector. In the $E_8$-Persistence framework, this is not an algebraic coincidence but a consequence of the single finite capacity of the lattice.

As established in Paper I (Section IX), the **Weak Mixing Angle** ($\sin^2 \theta_W$) is the partition ratio of the Temporal Resonance to the Total System Capacity:
\begin{equation}
\sin^2 \theta_W = \frac{\Delta}{N_{sys}} = \frac{43}{192.907} \approx \mathbf{0.2229}
\end{equation}

As derived in this work (Section III), the **Cabibbo Angle** ($\sin \theta_C$) is the leakage ratio of the Interface to the Active Flavor Width:
\begin{equation}
\sin \theta_C = \frac{\pi}{\nu - \chi} = \frac{\pi}{14} \approx \mathbf{0.2244}
\end{equation}

\subsubsection{Physical Interpretation: The Unified Partition}
We observe that the linear aperture of flavor mixing is structurally identical to the quadratic partition of force mixing:
\begin{equation}
\sin \theta_C \approx \sin^2 \theta_W
\end{equation}
\begin{equation}
|0.2244 - 0.2229| \approx 0.0015 \quad (< 0.7\%)
\end{equation}

This confirms that the Gauge Rotation ($\theta_W$) and the Mass Rotation ($\theta_C$) are two manifestations of the same \textbf{Geometric Partitioning Protocol}. The Selection Principle does not maintain two separate ledgers; it allocates the singular bandwidth ($H_{sys} \approx \Delta/2$) against the topological boundary cost ($\pi$ vs. $\Delta$) for both forces and masses simultaneously.
